% Glossar
\newglossaryentry{Entitaet}
	{name=Entität,
		description={Eine Entität, "ist ein Element der Datenwelt, welches ein reales oder gedankliches Einzelphänomen in einem betrachteten Realitätsausschnitt repräsentiert." (\cite{gehringWirtschaftsinformatik2022}: 787)},
	}

\newglossaryentry{Entitaets-Beziehung}
 	{name=Entitäts-Beziehung,
 		description={"Eine [Entitäts-]Beziehung (engl. relationship) ist eine logische Verknüpfung zwischen zwei oder mehreren Entitäten bzw. Entitätstypen." (\cite{gehringWirtschaftsinformatik2022}: 788)}
	}
	
\newglossaryentry{Entitaets-Typ}
	{name=Entitäts-Typ,
		description={"Ein Entitätstyp (oder Entitätsmenge, engl. entity set) fasst alle Entitäten zusammen, die durch gleiche Merkmale, nicht notwendigerweise aber durch gleiche Merkmalsausprägungen, charakterisiert werden." (\cite{gehringWirtschaftsinformatik2022}: 788)}
	}
	
\newglossaryentry{Assoziation}
{name=Assoziation,
	description={Im ERM Modell gibt eine "Assoziation [...] an, wie viele Entitäten eines Entitätstyps [A] einer beliebigen Entität des Entitätstyps [B] zugeordnet sein können." (\cite{gehringWirtschaftsinformatik2022}: 788)}
}

\newglossaryentry{Normalisierung}
{name=Normalisierung,
	description={}
}